% Transcription — fonds Alexandre Grothendieck, shelfmark 36, batch 4 (pages 61-80).
% Established from the facsimile archives/batches/36/batch-04.pdf (a copy of
% 36.pdf fetched through the project's relay, grothendieck-relay.onrender.com;
% PDF page k = archivists' page 60+k, no cover sheet), per the skill
% transcribe-grothendieck. The folder is 84 pages in five batches, done in
% parallel; this is the fourth.
%
% Pass: Opus 5.5 (claude-opus-5-5), 2026-09-26 — first pass, unchecked against the pages by a human.
%
% Pages skipped: 65, 67, 69, 71 (typed versos of an unrelated Bourbaki draft,
% « n° 384 », typed pp. 9, 7, 3, 10, on Haar measures of closed subgroups and
% the space of discrete subgroups, scanned upside down; nothing of SGA 7 on
% them) and 77 (a blank sheet of IHES letterhead). Page 72 is a cover (a sheet
% bearing only his title): no \page marker, his words become the section
% heading, with a \note.
% Pages summarised / noted only: none.
% Pages transcribed: 61-63 (typescript), 64, 66, 68, 70, 73-76, 78-80 (his hand).
%
% Typescript: pages 61-63 are the continuation of the typescript begun on
% p. 60 (batch 3) under his cover « "Th. du cycle invariant" : Cas du H^1 »
% (p. 59). transcripts/36/batch-03.fr.tex exists and decided it is an
% UNPUBLISHED typescript of his, transcribed in full; aligned with that
% decision. Evidence on these pages: the numbering continues batch 3's
% integer scheme (c), d), formulas (6)-(8), Corollaire 9 with (10),
% Corollaire 10 with (11)-(14), Corollaire 15 with (16)-(17)), not the decimal
% numbering of SGA 7; a search of a transcription of the published SGA 7 I
% (Exposé IX, §§ 11-12, where d, d^t and Raynaud's theorem 12.1 occur) finds
% none of these statements (no H^1(X, mu_n) -> H^1(X_etabar, mu_n)^I, no
% « Ker w_n », no Kummer covering of group mu_a); the ink and pencil
% corrections are an author's (« fonctorielle en X », « a fortiori si l ∤
% d^t », the substitution of Q(n) and of hooked arrows in (11)-(13)); and the
% handwritten pages 64-70 that follow are his own working of the same theorem
% in the same notation (Gamma, Delta, delta, delta', w_n, d'), with « demander
% à Raynaud si Ker est tout Z/dZ » on p. 70. The typescript breaks off in the
% middle of p. 63, after the statement of Corollaire 15; no page of it is
% numbered. Machine x-outs are not recorded unless they read; the pencil
% crosses in the left margin are not recorded.
%
% His pagination and numbered runs: none on these leaves. Typescript runs:
% a)-d), (6)-(17), Corollaires 9, 10, 15. Manuscript: « Théorème », « Dém. »,
% (*), (**) on p. 64; a circled (D) on p. 68; on pp. 73-80 « Théorème »,
% « Lemme », « Corollaire 1 », « Corollaire 2 », « Remarque », « Théorème ».
%
% What the batch is about. (a) Invariant cycles for H^1 (pp. 61-70): for
% X proper regular over a strictly local trait S, X_0 = sum d_i D_i,
% Gamma = Delta / Z.X_0, w : Gamma -> NS(X_0): exact sequences
% 0 -> _nGamma -> H^1(X, mu_n) -> H^1(X_etabar, mu_n)^I -> Ker w_n -> 0 and
% 0 -> H^1(X, Z_l(1)) -> H^1(X_etabar, Z_l(1))^I -> R (x) Z_l -> 0, R = Ker w
% cyclic of order d'; Cor. 9 (u(l) injective with finite cokernel, zero for
% almost all l), Cor. 10 (general n, with Theta(n) inside the n-torsion of
% Br(K), killed by p^m), Cor. 15 (after a tame Kummer base change of group
% mu_a the l-adic map becomes an isomorphism). The manuscript pages prove it
% via Kummer theory, Pic(X) -> Pic(X_eta) with kernel Gamma, and a lemma
% Pic(X)_n -> Pic(X_0)_n injective, citing Raynaud's note on the
% specialisation of Picard functors; p. 70 drops the 0-cycle of degree 1
% (Br(K) of p-torsion). (b) « Th. d'orthogonalité pour les modèles de Néron,
% démonstration arithmétique » (cover p. 72; pp. 73-80): for G, G' flat
% commutative groups over an adic regular complete X, abelian over U = X - Y,
% S of finite l-torsion and T' a torus over Y, the pairing phi_xi kills
% T_l(S~)|U x T_l(T~)|U; reduction to a trait with finite residue field and
% a lemma (over a finite field, any pairing T_l(S) x T_l(T) -> T_l(G_m) with T
% of multiplicative type vanishes, by Frobenius weights: q^2 vs q, q^{3/2} vs
% q, Weil); reduction of the general case by writing V as a limit of rings of
% finite type over Z (p. 76, « CQFD »). Corollaries for a Néron model G of A
% over a henselian DVR: M(A)^I_0 ⊂ M(A)^I ⊂ M(A) = T_l(A(Kbar)),
% phi(M^I, M'^I_0) = 0, and for a polarisation phi(M^I, M^I_0) = 0; a remark
% that the finite-level pairings need not vanish (pp. 78-79); p. 80, struck
% through, restates the theorem with both orthogonality relations and begins
% the limit argument again.
%
% Notation fixed once: \ell for the typed chi and for his l; \mathrm{NS} for
% NS; \underline{\mathbf{Z}}_\ell where the machine underlines, \mathbf{Z}_\ell
% in his hand; M^I_0 as he writes it (the torus part), M(A), M(A').
%
% Apparatus: 344 \ill{}, 212 \uncertain{} (the manuscript pages are fast
% drafting; the formulas carry, much of the connecting prose does not).
%
% Readings to check first: the inked letter before « -> S » on p. 63; the
% last line of p. 64; the struck four-line block and the marginal lemma of
% p. 66; most of the prose of p. 68 and p. 70 (fast drafting, heavily lost);
% the end of p. 73; the cover wording of p. 72.

\documentclass[11pt,a4paper]{article}
% Le préambule commun à toutes les transcriptions du fonds Grothendieck.
%
% Il tient en une idée : **l'incertitude se compose, elle ne se résout pas.**
% Une transcription qui lisse un mot douteux a détruit la seule chose pour
% laquelle on la faisait — un lecteur doit pouvoir savoir, à chaque endroit,
% ce qui était sur le feuillet et ce qui a été ajouté. Les six macros qui
% suivent sont donc l'appareil critique, et non de la décoration.
%
% Les mêmes commandes sont reconnues par `scripts/render.mjs`, qui produit la
% vue de lecture du site. Une macro ajoutée ici doit l'être aussi là-bas,
% faute de quoi le rendu échouera bruyamment — ce qui est voulu : un rendu
% silencieusement faux est pire qu'un rendu absent.

\ProvidesPackage{grothendieck}[2026/08/08 Transcriptions du fonds Grothendieck]

\usepackage{amsmath,amssymb,amsthm}
\usepackage{mathrsfs}
\usepackage{stmaryrd}
% Les diagrammes commutatifs sont partout dans ce fonds : c'est la langue
% même de la géométrie algébrique de Grothendieck, et une transcription qui
% les rendrait en prose ne transcrirait rien.
\usepackage{tikz-cd}
% Français par défaut — c'est la langue des feuillets — mais l'anglais est
% chargé aussi : la traduction et le résumé sont des documents anglais, et la
% typographie française (espaces avant les deux-points) y serait une faute.
% Ces documents basculent par \selectlanguage{english} en tête de corps.
\usepackage[english,french]{babel}
\usepackage{fontspec}
\usepackage{geometry}
\geometry{a4paper,margin=25mm,marginparwidth=28mm}
\usepackage{marginnote}
\usepackage{xcolor}
% ulem, not soul. `soul`'s \st and \ul are built on a character-by-character
% scan that XeLaTeX's font machinery breaks: the document fails with an
% undefined control sequence rather than degrading. ulem does the same two jobs
% and is Unicode-engine safe. `normalem` stops it hijacking \emph, which the
% transcriptions use for Grothendieck's own emphasis.
\usepackage[normalem]{ulem}
\usepackage{hyperref}
\hypersetup{colorlinks=true,linkcolor=black,urlcolor=black,pdfborder={0 0 0}}

% --- Métadonnées du lot -----------------------------------------------------
% Reprises telles quelles par le script de rendu, qui les affiche en tête de
% la vue de lecture. Elles fixent ce que le fichier prétend couvrir : sans
% elles, une transcription flotte sans point d'attache dans un fonds de seize
% mille feuillets.

\newcommand{\folder}[1]{\gdef\grothFolder{#1}}
\newcommand{\batch}[1]{\gdef\grothBatch{#1}}
\newcommand{\leaves}[2]{\gdef\grothFirst{#1}\gdef\grothLast{#2}}
\newcommand{\foldertitle}[1]{\gdef\grothFoldertitle{#1}}
\gdef\grothFolder{?}\gdef\grothBatch{?}\gdef\grothFirst{?}\gdef\grothLast{?}\gdef\grothFoldertitle{}

% --- Le résumé --------------------------------------------------------------
%
% Il ouvre la lecture modernisée, sous le titre « Résumé », et il est composé
% en petit. Ce n'est pas un ornement : c'est le seul passage du document qui
% ne soit pas des mathématiques, et un lecteur doit pouvoir le reconnaître
% avant de l'avoir lu — puis le sauter s'il connaît déjà le sujet, ou s'y
% arrêter s'il ne le connaît pas. Le filet qui le suit marque la reprise du
% corps.

\newenvironment{resume}
  {\small\setlength{\parskip}{0.45em}}
  {\par\medskip\hrule\medskip}

% --- Mots-clés ---------------------------------------------------------------
%
% En anglais, en clôture du résumé de la lecture modernisée : le vocabulaire
% moderne sous lequel chercher ce que les feuillets construisent. Le manifeste
% du site (`npm run manifest`) les extrait pour étiqueter la cote entière —
% cette ligne est la source unique des tags, il n'y a pas de fichier à côté.
\newcommand{\keywords}[1]{\par\smallskip\noindent
  {\footnotesize\textsc{Keywords} --- \textit{#1}}\par}

% --- L'appareil critique ----------------------------------------------------

% Le numéro de feuillet, en marge. C'est l'ancre qui permet de régler un
% désaccord sur une formule en regardant *un* feuillet plutôt qu'une chemise
% de six cents. Le site s'en sert aussi pour tourner les pages du fac-similé
% quand on fait défiler la transcription.
\newcommand{\leaf}[1]{%
  \marginnote{\footnotesize\color{gray!70}\textsf{#1}}%
  \label{leaf:#1}\ignorespaces}

% Illisible : on ne devine pas. Un mot inventé qui se lit comme les autres est
% le pire résultat possible.
\newcommand{\ill}{\textcolor{red!60!black}{[\,\ldots\,]}}

% Lecture douteuse : le mot est donné, et signalé comme incertain.
\newcommand{\uncertain}[1]{\uline{#1}}

% Ajout de l'éditeur — une lettre manquante, un mot sous-entendu.
\newcommand{\add}[1]{\textcolor{blue!50!black}{[#1]}}

% Biffé par Grothendieck lui-même. Ce qu'il a rayé fait partie du document :
% une rature dit souvent où la pensée a changé de direction.
\newcommand{\struck}[1]{\sout{#1}}

% Note du transcripteur, sur l'état matériel du feuillet.
\newcommand{\note}[1]{\par\smallskip{\small\color{gray!60!black}\textit{[#1]}}\par\smallskip}

% Note marginale de Grothendieck, distincte du corps du texte.
\newcommand{\marginal}[1]{\par\smallskip\noindent\hspace{1em}%
  {\small\color{gray!60!black}#1}\par\smallskip}

% --- Titre ------------------------------------------------------------------

\AtBeginDocument{%
  \begin{center}
    {\large\bfseries Fonds Alexandre Grothendieck --- cote n\textsuperscript{o}~\grothFolder}\\[2pt]
    {\itshape\grothFoldertitle}\\[4pt]
    {\small feuillets \grothFirst--\grothLast\ (lot \grothBatch)}\\[2pt]
    {\footnotesize\color{gray!60!black}Transcription établie d'après le fac-similé numérisé
     par l'Université de Montpellier.\\
     Les lectures incertaines sont soulignées, les illisibles notées [\,\ldots\,],
     les ajouts entre crochets.}
  \end{center}
  \vspace{1em}}

\endinput


\folder{36}
\batch{4}
\pages{61}{80}
\dating{[vers 1967-1973]}
\watermark{Édition de démonstration}
\foldertitle{SGA 7 (ma part) : notes manuscrites (s.d.), tapuscrit annoté (s.d.)}

\begin{document}

\section*{« Th.\ du cycle invariant » : Cas du $H^1$ (suite)}

\note{suite du tapuscrit commencé à la page 60 (lot précédent), sous le titre
inscrit de sa main sur le feuillet de garde de la page 59 ; le titre de section
en est repris. Tapuscrit à ruban bleu, corrigé à l'encre et au crayon ; le
$\ell$ de la machine est un caractère en forme de $\chi$, les $\mu$, les
$\bar\eta$ en indice et plusieurs flèches sont encrés à la main. Les biffures
faites à la machine ne sont pas reprises, sauf quand elles se lisent ; les
petites croix au crayon de la marge gauche ne sont pas relevées. Pas de
numéro de page du tapuscrit.}

\page{61}
d'ordre \add{$d'$} divisant $d$, et a fortiori $d^t$.
\note{« $d'$ » est ajouté au crayon au-dessus de « ordre ».}

c)\quad Soit $n$ un entier, et supposons que $n$ est premier à $p$,
$d_t$ est premier à $p$, ou que $k$ est parfait (donc
algébriquement clos). On a alors une suite exacte canonique, fonctorielle
en $X$ :
\[
(6) \qquad 0 \to {}_n\Gamma \to H^1(X, \mu_n) \xrightarrow{\ u_n\ }
H^1(X_{\bar\eta}, \mu_n)^I \to \operatorname{Ker} w_n \longrightarrow 0 ,
\]
\note{« $d_t$ est premier à $p$, ou que » est tapé dans l'interligne et
relié par un trait de crayon ; le « ou » de la ligne tapée est souligné, et
« $k$ est parfait (donc algébriquement clos) » est enfermé par un crochet de crayon. Sous ${}_n\Gamma$, au crayon : « $\simeq \mathbf{Z}/(n,d)\mathbf{Z}$ ».}
où
\[
w_n \colon \Gamma_n \longrightarrow \mathrm{NS}(X_0)_n
\]
est l'homomorphisme déduit de (5) par tensorisation par
$\underline{\mathbf{Z}}/n\underline{\mathbf{Z}}$, et où $u_n$ est
l'homomorphisme canonique (\quad). Pour $n$ variable multiplicativement, les
suites exactes (6) forment un système projectif de suites exactes, pour les
morphismes de transition évidents.

d)\quad Soit $\ell$ un nombre premier. Alors on \struck{trouve} une suite exacte canonique, \add{fonctorielle en $X$}
\struck{(obtenue, si $\ell$ satisfait \add{aux conditions de c)}, par passage à
la limite projective sur les puissances de $\ell$ dans les suites exactes
(6))} :
\[
(7) \qquad 0 \to H^1(X, \underline{\mathbf{Z}}_\ell(1))
\xrightarrow{\ u(\ell)\ } H^1(X_{\bar\eta}, \underline{\mathbf{Z}}_\ell(1))^I
\to R \otimes \underline{\mathbf{Z}}_\ell \to 0 ,
\]
où $R = \operatorname{Ker} w \simeq \underline{\mathbf{Z}}/d'\underline{\mathbf{Z}}$, donc
\[
(8) \qquad R \otimes \underline{\mathbf{Z}}_\ell \simeq
\begin{cases}
0 & \text{si } \ell \nmid d', \\
\underline{\mathbf{Z}}/\ell\underline{\mathbf{Z}} & \text{si } \ell \mid d' .
\end{cases}
\]
\marginal{en regard du premier cas, à l'encre : « a fortiori si $\ell \nmid d^t$, »}
\note{« fonctorielle en $X$ » est écrit au crayon au-dessus de la parenthèse,
que le crayon enferme dans un cadre barré de traits obliques ; « aux conditions
de c) » est ajouté au crayon dans la parenthèse. Au-dessus de « Alors », une
incise tapée puis effacée (« si $d_t$ premier à $p$ ou $k$ parfait ») se lit
encore. Dans la marge gauche, en regard de (8), un $\chi$ entouré au crayon, et
le numéro (8) est encadré au crayon. Le « $\ne$ » tapé dans « si $\ell \ne d'$ »
est surchargé d'un trait vertical : on lit « $\ell \nmid d'$ ».}

On obtient en particulier grâce à d) \struck{au moins si $\ell \neq p$} :

\emph{Corollaire} 9.\quad Soit $\ell$ un nombre premier. Alors l'homomorphisme
\[
(10) \qquad u(\ell) \colon H^1(X, \underline{\mathbf{Z}}_\ell(1)) \longrightarrow
H^1(X_{\bar\eta}, \underline{\mathbf{Z}}_\ell(1))^I
\]
est injectif, son conoyau est un groupe fini, nul pour presque tout $\ell$.
La même conclusion est vraie en remplaçant $\underline{\mathbf{Z}}_\ell(1)$ par
$\underline{\mathbf{Z}}_\ell$, à condition de supposer $\ell \neq p$.
\note{les numéros « 9 » et « (10) » sont écrits à la main ; après « premier, »
une restriction tapée (« si $\ell \ne p$ si $k$ est pas parfait \ldots ») est
effacée à la machine.}

\page{62}
Notons d'ailleurs, pour le cas d'un entier $n$ ne satisfaisant
\add{nécessairement} pas les conditions envisagées dans c) :
\note{« nécessairement » est tapé dans l'interligne.}

\emph{Corollaire} 10.\quad Soit $n$ un entier. On a alors une suite exacte
canonique
\[
(11) \qquad 0 \to {}_n\Gamma \to H^1(X, \mu_n) \xrightarrow{\ u_n\ }
H^1(X_{\bar\eta}, \mu_n)^I \to Q(n) \to 0 ,
\]
où $Q_n = \operatorname{Coker} u_n$ s'insère dans une suite exacte canonique
\[
(12) \qquad 0 \to \operatorname{Ker} w_n \to Q(n) \to \Theta(n) \longrightarrow 0 ,
\]
avec
\[
(13) \qquad \Theta(n) \hookrightarrow {}_n\operatorname{Ker}\bigl(\mathrm{Br}(K)
\to \mathrm{Br}(X_{\bar\eta})\bigr) \hookrightarrow {}_{(n,d^t)}\mathrm{Br}(K) ,
\]
de sorte que l'on a $\Theta(n) = 0$ lorsque $n$ ou $d^t$ sont premiers à $p$
ou que $k$ est parfait (donc $\mathrm{Br}(K) = 0$). En tous cas, si $p^m$ est
la plus grande puissance de $p$ divisant $d^t$ (qui est donc indépendant du
choix de $n$), on a
\[
(14) \qquad p^m \Theta(n) = 0 \qquad \text{pour tout } n .
\]
Pour $n$ variable multiplicativement, les diagrammes (11), (12) et (13)
forment des systèmes projectifs, pour des morphismes de transition évidents.
\note{« 10 » et les numéros (11) à (14) sont corrigés à l'encre sur des
numéros tapés (« 1\ldots »). Dans (11), le « $(n)$ » de $Q(n)$ est surchargé à
l'encre sur un indice. Dans (13), la première flèche est surchargée à l'encre
en « $\hookrightarrow$ », et le même signe est répété, entouré, dans la marge
gauche.}

Si donc $n$ parcourt les puissances d'un nombre premier fixé $\ell$, alors il
résulte de (14) que $\varprojlim_\nu \Theta(\ell^\nu) = 0$, donc
$\varprojlim_\nu \operatorname{Ker} w_{\ell^\nu} \simeq \varprojlim Q(\ell^\nu)$
par (12), d'où la suite exacte (7) par passage à la limite à partir de (11).
\note{l'indice $\ell^\nu$ de $w$ et le signe $\simeq$ sont ajoutés à l'encre ;
deux passages tapés sont effacés à la machine. Un double trait vertical à
l'encre, dans la marge droite, borde ce paragraphe.}

Nous prouverons également le

\emph{Corollaire} 15.\quad Posons
\[
(16) \qquad d^t = p^r a , \qquad \text{avec } (a,p) = 1 .
\]
Alors il existe un diagramme

\begin{tikzcd}
X \arrow[d, "f"'] & X' \arrow[l, "h"'] \arrow[d, "f'"] \\
S & S' \arrow[l, "g"']
\end{tikzcd}

(17), avec $h$ un revêtement étale \struck{principal de groupe}
\ill{}, $S' \to S$ un revêtement
\note{le numéro (17), les lettres $f$, $f'$, $g$ et les flèches verticales du
diagramme sont à la main ; « revêtement » est effacé à la machine devant
« diagramme ». Après le passage biffé, un signe à l'encre illisible.}

\page{63}
kummérien \add{connexe} (ramifié) de \add{$S$, de} groupe de Galois
$\mu_a$, le diagramme induit par le changement de base
\uncertain{$\eta$} $\to S$ étant cartésien. Un tel diagramme est unique à
isomorphisme près, et $f' \colon X' \to S'$ satisfait aux mêmes hypothèses que
$f$ dans \quad , mais l'entier $d(X'/S')$ (défini comme $d = d(X/S)$ dans (2))
est maintenant égal à $p^r$. Si $I' \subset I$ est le groupe d'inertie relatif
à $S'$, alors pour tout entier $n$ premier à $p$, l'homomorphisme
\[
u'_n \colon H^1(X', \mu_n) \longrightarrow H^1(X'_{\bar\eta}, \mu_n)^{I'}
\]
est injectif, et son conoyau est isomorphe à ${}_n\delta/{}_n\delta'$, où
$\delta$ est le sous-groupe de torsion de $\operatorname{Coker} w$ ($w$ défini
dans (5)) et $\delta'$ le sous-groupe de torsion de $\mathrm{NS}(X_0)$. On en conclut,
pour tout nombre premier $\ell \neq p$, que l'homomorphisme
\[
u'(\ell) \colon H^1(X', \underline{\mathbf{Z}}_\ell) \longrightarrow
H^1(X'_{\bar\eta}, \underline{\mathbf{Z}}_\ell)^{I'}
\]
est un \emph{isomorphisme}.
\note{« connexe » et « $S$, de » sont ajoutés à l'encre au-dessus de la ligne.
La lettre devant « $\to S$ » est tracée à la main sur un blanc de la machine ;
elle a la forme d'un $\eta$ (lecture douteuse, on attendrait $S' \to S$). Les
$\delta$, $\delta'$ et les $\mu$ sont à la main ; après les deux
$\underline{\mathbf{Z}}_\ell$ de $u'(\ell)$, un « (1) » tapé est effacé. « isomorphisme » est
souligné au crayon ; une suite tapée (« et la même conclusion \ldots ») est
effacée à la machine. Le tapuscrit s'arrête ici, au milieu du feuillet, sur
l'énoncé du Corollaire 15 ; aucune démonstration ne suit dans le dossier.}

\page{64}
\marginal{en haut à droite, encadré : « $\overset{\mathrm{dfn}}{=}
\operatorname{Im}\bigl(\mathrm{Pic}(X_0) \to
\underline{NS}_{X_0/k}(k)\bigr)$ »}

\emph{Théorème}\quad $S$ trait str.\ local,

$X$ schéma propre sur $S$, régulier, tel que $\underline{\mathcal{O}}_S
\xrightarrow{\sim} f_*(\underline{\mathcal{O}}_X)$, et que $X_\eta$ admette un
\uncertain{$0$-cycle} de degré $1$,

$\Delta$ le groupe des diviseurs sur $X$ tels que \uncertain{supp}
$\subset X_0$,

$c \colon \Delta \to \struck{\ill{}} = \underline{NS}_{X_0/k}(k)$ défini comme
composé
\[
\Delta \to \mathrm{Pic}(X) \to \mathrm{Pic}(X_0) \to
\underline{\mathrm{Pic}}_{X_0/k}(k) \to \underline{NS}_{X_0/k}(k) \to N
\]
\note{dans « $c \colon \Delta \to$ », un premier but biffé (« $NS(X_0)$ » ?),
puis « $\underline{NS}(X_0)$ » surchargé en $\underline{NS}_{X_0/k}(k)$ ; le
« $N$ » final de la ligne composée est tel qu'écrit.}

$\delta = $ \struck{\ill{}} sous-groupe de torsion de $\operatorname{Coker} c$,
\struck{\ill{} sous-groupe de torsion de $\underline{NS}(X_0)$, de sorte que
$\Delta \to \Lambda \xrightarrow{\lambda} \delta \to 0$}

Alors $\delta$ est un groupe \uncertain{fini} ; \struck{\ill{}}
\uncertain{torsion} $NS(X_0) = \delta'$, \uncertain{soit} $\delta'' =
\delta/\delta'$ ; et pour tout entier $n$ premier à $p$, on a une suite
exacte canonique
\[
0 \to {}_n\Gamma \to H^1(X, \mu_n) \xrightarrow{\ u_n\ }
H^1(X_{\bar\eta}, \mu_n)^I \longrightarrow {}_n\delta''
\to \struck{\ill{}}\ {}_n\delta' \to \delta_n \to \delta''_n \to 0
\]
\note{au-dessus de la flèche vers ${}_n\delta''$, une variante : « $\nearrow
{}_n\delta/{}_n\delta'$ » ; au bout, un « $\to \delta''$ » raturé. La lecture
de la fin de la suite est douteuse.}
\uncertain{compatible} avec les \ill{} de transition évidents \uncertain{en
$n$} ; \uncertain{d'où} par passage à la limite, pour tout $\ell$ nb.\ pr.\
$\neq p$
\[
H^1(X, \mathbf{Z}_\ell) \xrightarrow{\ \sim\ } H^1(X_{\bar\eta}, \mathbf{Z}_\ell)^I
\]
(\emph{Théorème des cycles invariants}).

\emph{Dém.}\quad $X$ est plat sur $S$, \uncertain{donc} $X_\eta$ dense dans
$X$, et $H^1(X, \mu_n) \hookrightarrow H^1(X_\eta, \mu_n)$.
\marginal{à droite, au-dessus : « (car $X$ \uncertain{intègre} »}
\uncertain{On a} le diagramme

\begin{tikzcd}[column sep=small, row sep=small]
H^1(X, \mu_n) \arrow[r, "\alpha"] \arrow[d] & {}_n\mathrm{Pic}(X) \arrow[d] & & \\
H^1(X_\eta, \mu_n)/(\mathbf{Z}/n\mathbf{Z}) \arrow[r, "\alpha_\eta"] \arrow[d] & {}_n\mathrm{Pic}(X_\eta) \arrow[r, "\beta_\eta"] \arrow[d] & {}_nP(K) \arrow[r, "\simeq"] \arrow[d] & {}_nP(\bar K)^I \arrow[dl, bend left=20] \\
H^1(X_{\bar\eta}, \mu_n) \arrow[r, "\alpha_{\bar\eta}"] & {}_n\mathrm{Pic}(X_{\bar\eta}) \arrow[r, "\beta_{\bar\eta}"] & {}_nP(\bar K) &
\end{tikzcd}

\note{les flèches $\alpha$, $\alpha_\eta$, $\alpha_{\bar\eta}$, $\beta_\eta$,
$\beta_{\bar\eta}$ portent chacune le signe $\sim$ ; l'arc de droite va de
${}_nP(\bar K)^I$ vers ${}_nP(\bar K)$.}

où $\alpha, \alpha_\eta, \alpha_{\bar\eta}$ sont des isom.\ par \uncertain{la
théorie} de Kummer et $f_*(\mathcal{O}_X) \simeq \underline{\mathcal{O}}_S$
(NB $V^*/V^{*n} = 0$,
\uncertain{$K^*/K^{*n}$} $= \mathbf{Z}/n\mathbf{Z}$, car $V$ str.\ loc.\@) ;
$\beta_\eta$ (et $\beta_{\bar\eta}$) \ill{} isom.\ à cause de l'existence d'un
$0$-cycle de degré $1$ sur $X_\eta$, \uncertain{donc} \ill{} \ill{} \uncertain{que}
\ill{} (*) s'identifie à :
\[
(**) \qquad u' \colon {}_n\mathrm{Pic}(X) \to {}_n\mathrm{Pic}(X_\eta)
\]
Comme $X$ est régulier (il suffirait : anneaux locaux factoriels), \ldots
\[
0 \to \Gamma \xrightarrow{\ i\ } \mathrm{Pic}(X) \xrightarrow{\ v\ }
\mathrm{Pic}(X_\eta) \to 0
\]
où $\Gamma = \Delta/\mathbf{Z}.X_0$, d'où
\[
\operatorname{Coker} u_n \simeq \operatorname{Coker} {}_nv \simeq
\operatorname{Ker} i_n , \quad \text{où } i_n \colon \Gamma_n \to
\mathrm{Pic}(X)_n .
\]
\note{le feuillet s'arrête ici. Les mots de liaison de la ligne suivant (**)
sont illisibles.}

\page{66}
\struck{Comme} \struck{On sait que} \struck{On sait que}

On a d'autre part, si
\[
j_X \colon \Gamma \to \mathrm{Pic}(X_0)
\]
est le composé, d'où des $j_n \colon \Gamma_n \to \mathrm{Pic}(X_0)_n$,
\struck{\ill{}} \uncertain{l'inclusion} évidente :
\[
\operatorname{Ker} i_n \subset \operatorname{Ker} j_n .
\]

\marginal{en regard, souligné : « \emph{Lemme} $\mathrm{Pic}(X)_n
\hookrightarrow \mathrm{Pic}(X_0)_n$ (du moins si $S$ complet, ce qui suffit
(\ill{}.)) \uncertain{facile par les} \ill{} »}
Cette inclusion est \add{en fait} une \emph{identité}. Pour le voir, on note
que si $L$ est un faisceau inv.\ sur $X$, et si $L_n = L|X_n$ est \ill{}
comme faisceau inversible, \uncertain{alors} l'obstruction \ill{} \ill{}
$L_{n+1}$ comme faisceau inversible est nulle (cf.\ PCA, exposé sur Picard).
Donc \uncertain{d'autre} part
\[
\mathrm{Pic}(X_0) = \underline{\mathrm{Pic}}_{X_0/k}(k)
\]
\uncertain{puisque} \ill{} $X_0$ \ill{} \uncertain{un cycle de degré $1$}, d'où
\[
0 \to \mathrm{Pic}^0(X_0) \to \mathrm{Pic}(X_0) \to NS(X_0) \to 0 ,
\]
avec $\mathrm{Pic}^0(X_0) \overset{\mathrm{dfn}}{=}
\underline{\mathrm{Pic}}^0_{X_0/k}(k)$ et $NS(X_0) \overset{\mathrm{dfn}}{=}
\operatorname{Im}\bigl(\underline{\mathrm{Pic}}_{X_0/k}(k) \to
\underline{NS}_{X_0/k}(k)\bigr)$, et $\mathrm{Pic}^0(X_0)$ est
\uncertain{$n$-divisible}, donc on a
\[
\mathrm{Pic}(X_0)_n \hookrightarrow NS(X_0)_n
\]
\marginal{à gauche, souligné : « \uncertain{prendre un} \ill{} \uncertain{principal} »}
donc si $k \colon \Gamma \to NS(X_0)$ est le composé $\Gamma \to
\mathrm{Pic}(X_0) \to NS(X_0)$, \ill{}
\[
\operatorname{Coker} u_n \simeq \operatorname{Ker} k_n .
\]

\marginal{à gauche, en regard : « cf.\ plus bas une démonstration »}
Je dis que \add{$\Gamma$ est sans torsion et que} $\delta$ est injectif i.e.\
qu'un faisceau inv.\ sur $X$ \uncertain{qui est} \ill{} sur $X_\eta$ est
trivial. \struck{\ill{} $X_\eta$, \ill{} $X_0$ \ill{} \ill{} \ill{} zéro, \ill{}
$NS(X_0)$ \ill{} \ill{} indique (\ill{}) que \ill{} \ill{} $\ell$-divisible pour
tout $\ell \neq p$, \ill{}}
\note{un long passage (quatre lignes) est encadré et barré de traits obliques ;
l'addition « $\Gamma$ est sans torsion et que » est écrite au-dessus de la
ligne, dans un cadre ; le « $\delta$ » est peut-être un $i$.}

\uncertain{On a} une suite exacte

\[
0 \to \Gamma \xrightarrow{\ k'\ } NS(X_0) \xrightarrow{\ \lambda\ } \delta \to 0
\]
\note{à gauche, une première suite « $0 \to \Gamma \xrightarrow{k} NS(X_0) \to
\ldots \to 0$ », dont le troisième terme est raturé, avec en dessous
« \ill{} » biffé ; au-dessus, « $\operatorname{Coker} k = \operatorname{Coker}
\Delta$ ». Sous $\delta$, à droite : « $=$ sous-gr.\ des $NS(X_0)$ dont un
multiple \uncertain{provient} de $\Gamma$ ».}
d'où \uncertain{par} \ill{} \uncertain{suite exacte} \ill{}
\[
\operatorname{Ker} k_n = \operatorname{Ker} k'_n \simeq \operatorname{Coker} {}_n\lambda .
\]
Comme $\Gamma$ est sans torsion, \ill{} $\delta' \xrightarrow{\lambda} \delta$
est injectif, et on trouve que $\operatorname{Coker} {}_n\lambda =
{}_n\delta/{}_n\delta'$. \ill{} \ill{} \ill{} \ill{}
\note{la dernière ligne, au pied du feuillet, n'est pas lue.}

\page{68}
exacte
\[
0 \to \delta' \to \delta \to \delta'' \to 0
\]
et la suite exacte \uncertain{associée}
\[
0 \to {}_n\delta' \to {}_n\delta \to {}_n\delta'' \to \delta'_n \to \delta_n
\to \delta''_n
\]

\emph{Reste} : prouver \ill{} le fait que $\Gamma$ est sans torsion. Mais
$\Gamma = \mathbf{Z}/d\mathbf{Z}$, \ill{} $d$ est le pgcd des $d_i$,
$X_0 = \sum d_i D_i$, $\ell X_{0,\mathrm{red}} = \sum D_i$.
\struck{\ill{} \ill{} si $\exists$ \ill{} $X_\eta$ un $0$-cycle de degré 1}
\uncertain{si} \ill{} \ill{} un pt fermé $x$ de $X_\eta$, \ill{} \ill{}
\add{\uncertain{schématique} dans $X$} est un \ill{} schéma fermé $Z$ de $X$,
plat et fini sur $S$, \struck{\ill{}} dont le degré sur $S$ est aussi égal à la
multiplicité d'intersection \add{$Z . X_0$} sur $X$, qui est
$\sum d_i\, Z . D_i$, donc multiple de $d$. Donc si \struck{\ill{}} $d'$ est le
pgcd des degrés des $0$-cycles sur $X_\eta$, \ill{} \ill{} $d \mid d'$. Comme
$d' = 1$, \ill{} $d = 1$.
\marginal{à gauche : « \uncertain{Lemme} (\uncertain{ici} $X$ \uncertain{régulier},
semble-t-il) »}
(D)\quad $\Gamma \to NS(X_0)$ est \emph{injectif}.\note{le « D » est
entouré.}
\marginal{à droite : « Soit $R$ \ill{} »}
\struck{\ill{}} \ill{} \uncertain{qui} \uncertain{utilise} \ill{} \ill{}
\uncertain{sans} \ill{}, \ill{} les suites \uncertain{exactes}
\marginal{à gauche, souligné en diagonale : « \uncertain{Donnons} \ill{}
\ill{}. »}
\[
0 \to H^1(X, \mu_n) \to H^1(X_{\bar\eta}, \mu_n)^I \to {}_n\Gamma_n \to
NS(X_0)_n
\]
\uncertain{d'où}, \uncertain{par passage} à la limite sur les puissances
\ill{} de $\ell$
\[
0 \to H^1(X, \mathbf{Z}_\ell(1)) \to H^1(X_{\bar\eta}, \mathbf{Z}_\ell(1))^I
\to T_\ell(\Gamma \otimes \ldots) \to T_\ell(NS(X_0) \otimes \mathbf{Z}_\ell)
\]
i.e.\ \ill{} \ill{}, \uncertain{comme} \struck{$T_\ell(\ldots)$}
$R \otimes \mathbf{Z}_\ell \simeq \operatorname{Ker}\bigl(\Gamma \otimes
\mathbf{Z}_\ell \to NS(X_0) \otimes \mathbf{Z}_\ell\bigr)$
\[
0 \to H^1(X, \mathbf{Z}_\ell(1)) \to H^1(X_{\bar\eta}, \mathbf{Z}_\ell(1))^I
\to R \otimes \mathbf{Z}_\ell \to 0
\]
\note{au-dessus de cette suite, « $T_\ell(R)$ » biffé, et un « $\beta$ » ;
les objets écrits $\Gamma$ et $\mathbf{Z}_\ell$ sont peu lisibles.}
\marginal{à gauche : « [ ] \uncertain{en remarquant} \uncertain{qu'ils} \ill{} »}
\uncertain{Démonstration} \ill{} \uncertain{invariante}, \uncertain{signifiant}
\ill{} \ill{} \ill{}. [\ill{} \ill{} \uncertain{invariants} ; \ill{} \ill{}
\ill{} de \uncertain{formule} des \uncertain{cycles} \ldots \ill{} : torsion
\uncertain{près}.]

\struck{L'idée est de se ramener au cas où $X$ est de dim.\ \ldots\ relative
\uncertain{un}. Dans ce cas, \ill{} \uncertain{sait} que $\Gamma$ est \ill{}
\uncertain{dit} \ill{} $(D_i . D_j)$ \ill{} \ill{} \ill{} $r-1$), \ill{} (\ill{}
le fait que \ill{} \ill{} \ill{} \ill{} $D_i$ \ill{}).}
\note{passage encadré et barré de traits obliques.}

Or c'est fait par Raynaud dans sa \uncertain{note} : spécialisation des
foncteurs de Picard, \emph{th.~3} \ill{} \ill{} \uncertain{sans} \ill{} \ill{}
\uncertain{projective}. \struck{\ill{}}
\marginal{à gauche, entouré : « \uncertain{démontrer} \uncertain{dév.}
\uncertain{suivants} \uncertain{au cas} \uncertain{de} \uncertain{dim.\ rel.}
1 »}
\note{les deux dernières lignes du feuillet sont biffées de traits ondulés ;
on y lit $NS(X_0)$, sans plus.}

\page{70}
\emph{Variante}\quad Laissons tomber l'hypothèse $d'$ ($=$ pgcd des deg.\ des
$0$-cycles sur $X_\eta$) $= 1$.
On sait que $\mathrm{Br}(K)$ est un groupe de $p$-torsion, donc \ldots
\[
\to \operatorname{Pic}(X_\eta) \to P(K) \to \Theta \to 0 , \quad \text{où }
\Theta \subset \mathrm{Br}(K)
\]
est un groupe de $p$-torsion, d'où \uncertain{résulte} que
${}_n\operatorname{Pic}(X_\eta) \simeq {}_nP(K)$. \ill{} \ill{} \ill{} \ill{} que
\ill{} est isom.\ \uncertain{si} \ill{}. Mais $\Gamma$ peut \uncertain{avoir} de
la torsion. \ill{} \uncertain{trouve}
\[
0 \to {}_n\Gamma \to H^1(X, \mu_n) \to H^1(X_{\bar\eta}, \mu_n)^I \to
\Gamma_n \to \mathrm{NS}(X_0)_n \to
\]
\note{sous ${}_n\Gamma$ : « $\simeq \mathbf{Z}/(n,d)\mathbf{Z}$ ».}
\struck{Soit} et \uncertain{par} passage à la limite sur les puissances de
$\ell$
\[
0 \to H^1(X, \mathbf{Z}_\ell(1)) \to H^1(X_{\bar\eta}, \mathbf{Z}_\ell(1))^I
\to \Gamma \otimes \mathbf{Z}_\ell \to \mathrm{NS}(X_0) \otimes \mathbf{Z}_\ell
\]
\[
\hookrightarrow R \otimes \mathbf{Z}_\ell \to 0 , \qquad R = \operatorname{Ker}
\bigl(\Gamma \to \mathrm{NS}(X_0)\bigr)
\]
\marginal{à gauche, encadré : « N.B. $R \otimes \mathbf{Z}_\ell \simeq
\mathbf{Z}/\ell\mathbf{Z}$ si $\ell \mid d$, $0$ sinon »}
On voit ainsi (\ill{} au moins modulo rés.\ des singularités, \uncertain{pour}
\ill{} \ill{} \ill{} \ill{} \ill{}) que le noyau $R$ de $\Gamma \to
\mathrm{NS}(X_0)$ est de \emph{torsion}. \struck{\ill{} \ill{} \ill{} \ill{}
\ill{} \ill{}}
\marginal{à gauche, sous un trait : « \uncertain{mieux} \ill{}
\uncertain{vérifier} \uncertain{Raynaud} \ill{} \ill{}, \uncertain{p.\ 7} (ii) »}
\struck{\ill{}} (C'est donc un sous-groupe du \uncertain{tors} $\Gamma \simeq
\mathbf{Z}/d\mathbf{Z}$, \ill{} \ill{} \uncertain{cette} \uncertain{tendance}
d'être égal à $\mathbf{Z}/d\mathbf{Z}$ \ill{} \ill{}, \ill{} \ill{} \ill{} il le
\uncertain{faudrait} \ill{} \uncertain{tors} $\Gamma$ \ill{} \ill{} \ill{} le
noyau \ill{} \ill{} Raynaud \uncertain{utilisé} \ill{} \ill{} : \ill{}
\uncertain{arguments de} Raynaud \ill{} \ill{} \ill{} \ill{} de $S$ \ldots)
\ill{} \ill{}, si
\marginal{à gauche : « \uncertain{demander} à Raynaud si Ker est
\uncertain{tout} $\mathbf{Z}/d\mathbf{Z}$. »}
\[
\Gamma' = \Gamma/\struck{\ill{}}\,T, \qquad T = \operatorname{Tors} \Gamma
\]
\[
0 \to T_n \to \Gamma_n \to \Gamma'_n \to 0
\]
\ill{}, \uncertain{utilisons} un splitting de $\Gamma \simeq \Gamma' \oplus T$,
\ill{} \ill{}
\[
\operatorname{Ker}\bigl(\Gamma_n \to \mathrm{NS}(X_0)_n\bigr) \simeq T_n +
\underbrace{\operatorname{Ker}\bigl(\Gamma'_n \to \mathrm{NS}(X_0)_n\bigr)}_{\operatorname{Ker}(\bar\Gamma_n
\to \mathrm{NS}(X_0)_n)}
\]
\ill{} $\bar\Gamma \hookrightarrow \mathrm{NS}(X_0)$, $\bar\Gamma = \Gamma/R$.
\ill{} \ill{} \ill{}, \uncertain{comme} ${}_n\bar\Gamma$ \ill{} \ill{} \ill{}
$p$-torsion
\[
0 \to \delta' \to \delta \to \delta'' \to 0
\]
\note{sous $\delta'$ : « $=$ tors.\ de $\mathrm{NS}(X_0)$ premier à $p$ ».
Le feuillet s'arrête ici.}

\section*{Th.\ d'orthogonalité pour les modèles de Néron, démonstration arithmétique}

\note{titre inscrit de sa main en haut à droite du feuillet de garde (page 72),
qui ne porte rien d'autre ; le feuillet ne reçoit pas de numéro de page. On lit :
« Th.\ d'orthogonalité \uncertain{$\ell$/p} pour les modèles de Néron,
démonstration arithmétique », et plus bas, du même crayon, « $\ell$
(partiel !) ». Le signe au-dessus de « orthogonalité » est peut-être un
$\ell$ repris ; lecture douteuse. Les pages 73 à 80 suivent sans pagination
de l'auteur.}

\page{73}
\emph{Un théorème d'orthogonalité}

\struck{\ill{}}
$A$ anneau noethérien $J$-adique (régulier et complet)

$X = \operatorname{Spec} A$

$Y = \operatorname{Spec}(A/J) = V(J)$

\struck{\ill{}} $U = X - Y$

Pour tout schéma fini étale $Z$ sur $Y$, $\tilde Z$ est le schéma fini étale
sur $X$ (unique à isom.\ unique près) le prolongeant. \add{\uncertain{Dit}
\ill{} \uncertain{séparés}} Donc \ill{} \ill{} faisceaux $\ell$-adiques
\uncertain{constants tordus} \ldots
\note{« $\tilde Z$ » avec tilde sur $Z$ ; l'ajout est interlinéaire, sur un
mot biffé.}

$G, G'$ schémas \add{en groupes commutatifs} \struck{\ill{}} \add{plats}
\struck{\uncertain{séparés}} sur $X$
\marginal{à droite, encadré : « $A$ et $A'$ \struck{accouplés}
\struck{\ill{}} i.e.\ on a une \uncertain{classe} de corr.\ div.\ sur
$A \times A'$ »}

$G_U = G|U$ est un schéma abélien $A$

$G'_U = G'|U$ \quad ------ \quad $A'$

$\ell$ \struck{\ill{}} est premier \uncertain{inv.}\ sur $Y$, donc sur $X$.
\marginal{à gauche, entouré : « \uncertain{Composantes} \ill{} \ill{}
\uncertain{de} $G$, \ill{} \uncertain{dans} \ill{} \ill{} $A'$, $A$ »}

$S \subset G_Y = G \times_X Y$ est \add{plat sur $Y$ et} tel que ${}_{\ell^\nu}S$
\uncertain{fini} sur $Y$ pour tout $\nu > 0$
\struck{\ill{}}

$T' \subset G'_Y$ un \emph{tore}
\struck{Alors \ldots}

\marginal{à gauche : « NB les ${}_{\ell^\nu}G$ sont étales et q.f.\ \uncertain{sép.}\
sur $S$, et leur restriction à $Y$ contient les ${}_{\ell^\nu}S$, donc
${}_{\ell^\nu}G$ contient ${}_{\ell^\nu}\tilde S$. »}
$T_\ell(\tilde S)$ est \uncertain{bien} défini, et il en est de même de
$T_\ell(\tilde T)$ \uncertain{et} \ldots
\[
T_\ell(\tilde S) \subset T_\ell(G), \qquad T_\ell(\tilde T) \subset T_\ell(G'),
\]
\[
T_\ell(\tilde S)|U \subset T_\ell(A), \qquad T_\ell(\tilde T)|U \subset T_\ell(A') .
\]
\[
\varphi_\xi \colon \ldots \longrightarrow T_\ell(\mathbb{G}_{m,U}) =
T_\ell(\widetilde{\mathbb{G}_{m,Y}})|U
\]
\note{la flèche verticale $\varphi_\xi$ est suivie de « \uncertain{pairing} :
$\xi$ induit » ; à droite, « $T_\ell(S)$ ». Mise en ligne de l'éditeur.}

\emph{Théorème}\quad $\varphi_\xi\bigl(T_\ell(\tilde S)|U, T_\ell(\tilde
T)|U\bigr) = 0$ \ill{} \ill{} \ill{} \uncertain{en chaque} pt $x \in U$.

\emph{Dém.}\quad Il suffit de le \uncertain{prouver} en \uncertain{chaque} pt
$y \in \overline{\{x\}} \cap Y$, \struck{\ill{}}, et \uncertain{donnons} \ill{}
\uncertain{Prenant} un point \ill{} \ill{} \uncertain{discret}, \ill{} \ill{}
\ill{} \ill{} \ill{} \ill{} \ill{} \ill{}. Je \ill{} $X$ \ill{} un \emph{trait},
$Y$ \ill{} pt fermé, $U$ \ill{} pt génér. Je \uncertain{suffit} d'ailleurs,
\ill{} \ill{} \uncertain{réduction} \ill{} \uncertain{formule} \ill{}
\uncertain{suffit} \ill{} \ill{} \ill{} \ill{} \ill{} $U$ : $S$-\ill{} \ill{}
\ill{} \ill{}
\note{les dernières lignes de la page, en écriture rapide, ne sont lues que
par bribes.}

\page{74}
\uncertain{puisque}, \uncertain{pour} tout tel $x$, \uncertain{choisir}
$y \in \overline{\{x\}} \cap Y$ tel que $k(y)$ soit \emph{fini}
\add{(\uncertain{c'est le cas si} $A/J$ \uncertain{est de type fini sur}
$\mathbf{Z}$)}. Alors \struck{par} l'argument \uncertain{précédent}, \ill{}
\uncertain{on est} ramené au cas \ill{} \add{\uncertain{localisé} complété de}
\struck{\ill{} \ill{} \ill{} \uncertain{discret}} $X$ en $y$,
\struck{\ill{} \ill{}} donc $X$ local complet, de pt fermé $y$, \uncertain{avec}
$k(y)$ fini. Alors $X - Y$ est de \uncertain{Jacobson}, \ill{} tout
\uncertain{point} \ill{} de $X - Y$ \ill{} un pt $z$ tel que $\dim
\overline{\{z\}} = 1$, $\overline{\{z\}} = \overline{\{z\}}$ \ill{}
\uncertain{l'adhérence} dans $X$. \ill{} \ill{} \ill{}, \struck{\ill{}}
\ill{} : \uncertain{remplaçons} $X$ par $Z$, \ill{} \ill{} \ill{} $X$ local
complet de dim.\ $1$, et \uncertain{normalisons}, \ill{} \ill{} \ill{} $X$ est
\supplied{le spectre d'}un anneau de valuation discrète \struck{\ill{}} à corps
résiduel \uncertain{fini}, $Y$ pt fermé. L'\uncertain{homomorphisme} envisagé
correspond \ill{} \ill{} \ill{} \ill{} $\pi = \operatorname{Gal}(\bar K/K)$,
\uncertain{restreint} à \ill{} \uncertain{groupe} d'inertie $I$.
\struck{de $\pi_0 = \operatorname{Gal}(\bar k/k)$ \ill{} \ill{}
$\hat{\mathbf{Z}}$ \ldots\ $T \simeq \mathbb{G}_{m,k}^r$ (\ill{} \ill{}
résiduelle finie) et $T_\ell(\tilde T)|U$ \uncertain{correspond} \ldots\
$T_\ell(\bar K^*) \simeq T_\ell(\bar k^*)$, $\bigl(T_\ell(\mathbb{G}_{m,k})|U\bigr)
\simeq T_\ell(\mathbb{G}_{m,k}) \simeq T_\ell(k^*)$, comme modules sur
$\pi_0$. D'autre part, \ill{} $S$ \ill{} $T_\ell(S)$ \uncertain{correspond}.}
\note{un passage de six lignes est encadré et barré de traits obliques.}
L'accouplement sur $U$, \uncertain{comme} $X$ \ill{}, \uncertain{provient}
d'un accouplement
\[
T_\ell(S) \times T_\ell(T) \to T_\ell(\mathbb{G}_{m,k}) .
\]
Donc \ill{}, \ill{} \ill{} est ramené au

\emph{Lemme}\quad Soit $k$ un corps fini, $S$, $T$ deux \struck{\ill{}}
\uncertain{groupes} alg.\ commutatifs sur $k$, \struck{\ill{} \ill{}}
\add{$T$ de t.\ mult.} \uncertain{tore}. Alors tout accouplement
\[
T_\ell(S) \times T_\ell(T) \to T_\ell(\mathbb{G}_m)
\]
est identiquement nul.

\page{75}
En effet, \struck{\ill{}} on \uncertain{peut} supposer \add{$S$, $T$
\uncertain{connexes} \uncertain{réduits}, donc $T$ un tore,}
\struck{\uncertain{réduits} ; \ill{} \ill{} extension \add{finie} de $k$,
\ill{} $T$ \ill{} un tore \ill{} $S$ \ill{} $S$}
On \uncertain{peut} supposer aussi $S$ extension d'une VA par un tore $S_0$.
Quitte à faire une extension finie, on peut supposer $S_0 \simeq
\mathbb{G}_m^r$, $T \simeq \mathbb{G}_m^s$. Alors
\[
T_\ell(S_0) \times T_\ell(T) \to T_\ell(\mathbb{G}_m)
\]
\uncertain{est} un accouplement $\mathbf{Z}_\ell^r \otimes \mathbf{Z}_\ell^s \to
\mathbf{Z}_\ell$, \uncertain{où} \ill{} \ill{}. \ill{} \uncertain{le} Frob.\
opère \uncertain{sur} \ill{} les trois facteurs \uncertain{par} $q\,\mathrm{id}$,
$q\,\mathrm{id}$, $q\,\mathrm{id}$, \uncertain{donc} \ill{} \ill{} \ill{}
\uncertain{produit} \ill{} \uncertain{par} $q^2\,\mathrm{id}$, \ill{} \ill{}
\ill{} \uncertain{par} $q\,\mathrm{id}$. Donc l'accouplement induit est nul.
\uncertain{Cela} \ill{}, \ill{} \ill{} \ill{} \ill{} \ill{} $S$ \uncertain{une}
V.A. Mais, par Weil, les val.\ propres de $f$ \uncertain{opérant} sur
\struck{\ill{}} $T_\ell(S)$ \uncertain{sont} de val.\ absolue $q^{1/2}$, donc
les val.\ propres sur $T_\ell(S) \otimes T_\ell(\mathbb{G}_m)$ sont de val.\
absolue $q^{3/2}$, \ill{} \ill{} \uncertain{sur} $T_\ell(\mathbb{G}_m)$ \ill{}
$q$, \ill{} \ill{} \ill{} \uncertain{nullité}.

Nous traitons maintenant le cas général, \struck{\ill{}} \ill{} \ill{}
\uncertain{supposons}, \ill{} \ill{} \ill{} \ill{} $X$ \uncertain{un}
\struck{\ill{}} \add{$= \operatorname{Spec} V$, \ill{} \uncertain{de
type fini}} \ill{} \ill{} \ill{} pt fermé. \ill{} \ill{}
\[
V = \varinjlim A_i , \quad A_i \subset V , \quad A_i \ni \ldots \text{ de t.\ f.\ sur } \mathbf{Z} .
\]
\struck{\ill{}} Soient $X_i = \operatorname{Spec} A_i$, $Y_i = V(J_i)$, $U_i =
X_i - Y_i$, donc
\[
X = \varprojlim X_i , \quad Y = \varprojlim Y_i , \quad U = \varprojlim U_i .
\]

\page{76}
Pour $i$ grand, $G$, $G'$ proviennent de $G_i$, $G'_i$ \add{sch.\ en groupes
commut.\ de t.f.}\ plats sur $X_i$, \struck{et $\xi$ d'une corr.}
$G_i|U_i = A_i$ et $G'_i|U_i = A'_i$ des sch.\ abéliens, $\xi$ provient d'une
classe de corr.\ div.\ $\xi'$ sur $A_i \times A'_i$, $T$ provient d'un tore
$T_i \subset G'_{Y_i}$, $S$ provient d'un sous-schéma plat $S_i \subset
G_{Y_i}$. De plus, \struck{\ill{}} \ill{} on peut \struck{\ill{}
\uncertain{conditions}} supposer $S$ connexe \add{(\uncertain{car} on
\uncertain{peut} le remplacer par $S^0$)} et \struck{\ill{} \ill{} \ill{}}
\struck{\ill{} \ill{} \ill{}} \uncertain{extension} \struck{\ill{}}
\uncertain{successive} d'une V.A., d'un tore \struck{\ill{} \ill{} \ill{}}
\uncertain{et d'un} \uncertain{groupe} \ill{}. On peut \uncertain{supposer} la
\uncertain{même} \ill{} \ill{} pour $S_i$. \struck{\ill{}} On peut supposer
aussi (\uncertain{remplaçant} au besoin \struck{$A_i$} $A$ par $A_i[1/\ell]$)
que $\ell$ \uncertain{inv.}\ sur les $X_i$, \ill{} \ill{} sur les $Y_i$. Alors
les ${}_{\ell^\nu}S_i$ sont \uncertain{connexes} \ill{} \ill{} \uncertain{finis}
sur $Y_i$. Bref, \uncertain{tous} les données \uncertain{analogues}, avec
\uncertain{les} \ill{} \uncertain{provenant} de \ill{} données \ill{} \ill{}
\add{$X_0 = \operatorname{Spec} A_0$} de type fini sur $\operatorname{Spec}
\mathbf{Z}$. De plus, si $\mathfrak{p} = \mathfrak{m} \cap A_0$, alors $A_0
\to V$ induit $\hat A_0 \to V$, \struck{\ill{}}. Donc, \ill{} \ill{}
\uncertain{remplaçant} \ill{} \ill{} « $A_0$ de type fini sur $\mathbf{Z}$ »
\uncertain{par} « $A_0$ $\mathfrak{p}$-adique régulier et complet, avec
$A_0/\mathfrak{p}$ de t.f.\ sur $\mathbf{Z}$, et $A_0 \to V$ continue pour les
topologies \uncertain{adiques} ». On \uncertain{est} ramené \uncertain{alors}
au \uncertain{théorème} pour \ill{} $X_0, Y_0, G_0, G'_0, \xi_0, S_0, T_0$, et
il \uncertain{relève} du cas 1.\quad CQFD.

\page{78}
\emph{Corollaire 1}\quad $V$ anneau de val.\ discrète \add{hensélien, corps
résiduel $k$}, $K$ corps des fractions, \add{$\bar K$ clôture alg.}, $A$
schéma abélien sur $K$, $G$ modèle de Néron de $A$, \add{$T$ tore maximal de
$G_k$,} \struck{\ill{} dans $T_\ell(A)$} $\ell$ nb premier $\neq$ car.\ $k$,
\struck{\ill{} $T_\ell(\tilde G_k)|U$ \ill{}}
\[
T_\ell(\tilde T)|U \subset T_\ell(\tilde G_k)|U \subset T_\ell(A)
\]
et en prenant les points à valeurs dans $\bar K$
\[
M(A)^I_0 \subset M(A)^I \subset M(A) = T_\ell(A(\bar K))
\]
\note{les deux premiers termes sont surchargés (un premier « $M(T)$ » corrigé
en $M(A)$), un « $T_\ell(A(\bar K))$ » final est biffé ; le $0$ en indice est
tel qu'écrit.}
où $I$ est le groupe d'inertie $\subset \pi = \operatorname{Gal}(\bar K/K)$,
et où l'exposant $I$ désigne les invariants sous $I$.
Ceci posé, si $A'$ est un \uncertain{schéma abélien} \add{sur $K$},
et $\xi$ une \struck{\ill{}} corr.\ divisorielle sur $A \times A'$, d'où un
accouplement compatible avec $\pi$
\[
\varphi_\xi \colon M(A) \times M(A') \to T_\ell(\bar K^*) \struck{= \ldots} ,
\]
\uncertain{d'où} des \emph{\uncertain{relations}} \uncertain{entre} $M(A)^I$ et
$M(A')^I_0$, \ill{} \uncertain{entre} $M(A)^I_0$ et $M(A')^I$.

\uncertain{généraux}, ainsi que \add{En particulier}

\emph{Corollaire 2}\quad Soit $\xi$ une \emph{polarisation} sur $A$, d'où
\uncertain{une} forme bil.\ alternée non dégénérée
\[
\varphi_\xi \colon M(A) \times M(A) \to T_\ell(\bar K^*)
\]
Alors
\[
\varphi_\xi(M^I, M^I_0) = 0 .
\]
En effet, \uncertain{les} \ill{} \uncertain{corollaires} \ill{} \ill{}
\uncertain{équivalents}, \ill{} \ill{} \ill{} \ill{} \ill{} \ldots

\emph{Remarque}\quad \add{Si $k$ est sép.\ clos,} à la limite projective,
\uncertain{les} \uncertain{accouplements}
\[
{}_{\ell^\nu}A'(K) \times {}_{\ell^\nu}T(k) \to {}_{\ell^\nu}k^* \simeq {}_{\ell^\nu}K^*
\]
sont nuls. Cependant ces accouplements \ill{} \ill{} \ill{}

\page{79}
\ill{} \uncertain{pas} \uncertain{nuls} \ill{} \ill{}, \ill{} \ill{}
\uncertain{même} pour $\nu = 1$. En effet, \uncertain{si} \struck{\ill{}}
$x \in {}_{\ell^\nu}T(k)$ \add{$\subset {}_{\ell^\nu}A(K) \subset
{}_{\ell^\nu}A(\bar K)$} \uncertain{est} \ill{} \uncertain{non nul}, il existe
un $x' \in {}_{\ell^\nu}A'(\bar K)$ tel que \struck{\ill{}} $\varphi(x, x')
\neq 0$, et \struck{\ill{}} \uncertain{quitte} à agrandir $K$ \uncertain{en}
$K'$, on \uncertain{aura} donc un accouplement ${}_{\ell^\nu}A'(K) \times
{}_{\ell^\nu}T(k) \to {}_{\ell^\nu}k^*$ \uncertain{non} nul \ldots
\marginal{à gauche, reliée par un trait : « \ill{} \uncertain{par}
$A'$ \uncertain{le dual de} $A$. »}
\note{le reste du feuillet est blanc.}

\page{80}
\note{la page entière est barrée de deux longs traits obliques ; elle se lit
néanmoins.}

\emph{Théorème}\quad $V$ anneau de val.\ discrète \add{str.\ hensélien},
corps fr.\ $K$, corps rés.\ $k$,

$A$ schéma abélien sur $K$, modèle de Néron $G$

$A'$ \quad ------ \quad dual, \quad ------ \quad $G'$

\struck{$A_K$} $\bar K/K$, $T_\ell(A(\bar K))$ \qquad $T$ tore maximal de
$G_k$, $T'$ \quad ------ \quad $G'_k$,
\[
\begin{aligned}
T_\ell(A(\bar K)) \supset T_\ell(A(\bar K))^I &= T_\ell\bigl(A(\bar K)^I\bigr)
= T_\ell(A(K)) = T_\ell(G(V)) \\
&= T_\ell(G(k)) = T_\ell(G^0(k)) \supset T_\ell(T(k))
\end{aligned}
\]
\note{sous le premier terme, « $\shortparallel$ $M$ » ; en fin de première
ligne, un terme raturé.}
\[
\begin{gathered}
M \supset M^I \supset M^I_0 \\
M' \supset M'^I \supset M'^I_0
\end{gathered}
\qquad \text{de } \ldots
\]
On a
\[
M \times M' \xrightarrow{\ \varphi\ } T_\ell(\mathbb{G}_m(\bar K)) =
T_\ell(\bar K^*) \quad \bigl[\ \overset{\text{non can.}}{\simeq}
\mathbf{Z}_\ell\ \bigr]
\]
Je dis que
\[
\begin{cases}
\varphi(M^I_0, M'^I) = 0 \\
\varphi(M^I, M'^I_0) = 0
\end{cases}
\]

\emph{Dém.}\quad Pour \uncertain{simplifier}, supposons ${}_\ell G_k \subset
G^0_k$. \ldots
\[
V = \varinjlim A_i , \quad A_i \subset V \text{ de t.f.\ sur } \mathbf{Z},
\ \text{\uncertain{avec} } t, \ \text{\uncertain{normaux}}.
\]
$G$ provient de $G_i$, $G'$ de $G'_i$. $T$ de $T_i$, $T'$ de $T'_i$.
\struck{On peut supposer que \ldots}
\[
X_i = \operatorname{Spec}(A_i), \quad Y_i = V(t_{X_i}), \quad U_i = X_i - Y_i
= X_{i,t} ,
\]
\[
{}_\ell G(k) = \varinjlim {}_\ell G_i(Y_i), \qquad {}_\ell G'(k) = \varinjlim
{}_\ell G'_i(Y_i)
\]
\marginal{à droite : « On peut supposer $Y_i$ \uncertain{connexe}, donc les
\uncertain{morphismes} de transition injectifs. Pour $i$ grand, ils sont
bijectifs. »}
\[
{}_\ell G(k) \simeq {}_\ell G_i(Y_i), \qquad {}_\ell G'(k) \simeq {}_\ell
G'_i(Y_i) \qquad \Big|\ {}_\ell T
\]
\note{le feuillet s'arrête ici ; « $t$ » est l'uniformisante de $V$.}

\end{document}
